% Ad Atticum 12.40 --- headnote
% ad-atticum-12-40, 9 May 45 BC, Astura

\textit{Cicero to Atticus, written from Astura on the
seventh day before the Ides of May 709 AUC --- 9 May
45 BC (the manuscript dateline: \textit{Scr.\ Asturae
vii Id.\ Mai.\ a.\ 709 (45)}). The longest and most
defensive letter of the early-May cluster. Hirtius
has sent Cicero a draft attack on Cato that flatters
Cicero himself with ``the greatest praises''; Caesar's
own \textit{Anticato} is in preparation, and Cicero
treats Hirtius' piece as a useful preview. He has the
book copied at once through Musca and Atticus'
copyists --- he wants it ``widely circulated.''}

\textit{The middle of the letter is the answer to a
question Atticus had raised: that Cicero's standing
(\textit{gratia} and \textit{auctoritas}) may be
suffering from his prolonged \textit{maeror}. Cicero
flares back: what is he supposed to do --- not
grieve? not lie low? At the same time he registers,
in a single line, the failed \textit{symbouleutikon},
the political-advisory tract addressed to Caesar in
the manner of Aristotle's and Theopompus' letters to
Alexander --- \textit{[Greek: sumbouleutikon]},
\textit{[Greek: Aristotelous]}, \textit{[Greek:
Theopompou pros Alexandron]}. He cannot find the
shape of it: those Greeks wrote what was honourable
to themselves and pleasing to Alexander, and the
parallel will not bear weight. The catalogue of his
recent productivity that follows --- thirty days at
the gardens, no visitor turned away, the people with
him finding leisure harder to bear than he the labour
--- is one of the strongest defences of his grief-work
in the whole sequence.}

\textit{Section 3 turns the question outward: he is
not in Rome because it is the court recess
(\textit{discessus}); not at his seasonal estates
because he could not bear the \textit{frequentia} ---
a different word from the \textit{celebritas} of the
\textit{fanum} project, but the same vocabulary of
crowds. The closing remark on the lost
\textit{hilaritas} ``with which I used to season this
bitter time'' is the most explicit acknowledgment
that the man who returns to Rome will not be the man
who left it. Sections 4--5 are pure business: the
Scapula gardens (still the favoured site for the
shrine) are to be put up by the auctioneer through
joint pressure; the daggered crux \dag\ \textit{non est
in eo}\dag\ over Lentulus is preserved.}
